Across nine models spanning 1.4B--14B parameters and seven architecture families (MHA through GQA 8:1), we find that per-head non-uniform KV cache quantization on hardware-aligned bit-widths $\{2,4,8,16\}$ provides negligible improvement over uniform 4-bit for eight of nine models.
The sole exception is Qwen2.5-7B, where a cardinality-constrained solver achieves 10.91 perplexity at $4\times$ compression---a 29.4\% reduction over uniform 4-bit (15.45)---while unconstrained solvers degrade catastrophically (32--37 PPL) due to 2-bit distortion compounding through the residual stream.
The benefit depends not on GQA ratio or model scale, but on whether sensitivity heterogeneity is \emph{spread} across multiple heads with cross-bitwidth gradients---a pattern observed in only one model ($N{=}1$) in our study.
On the remaining eight models---including Qwen2.5-14B, where scaling dilutes the spread pattern---non-uniform allocation provides no meaningful benefit: Yi-1.5-9B has high concentration (0.71, exceeding Qwen-7B's 0.61) yet gains nothing, because its single-outlier ``spike'' has no cross-bitwidth gradient.
The cardinality constraint prevents catastrophic over-allocation on Qwen-7B and moderate degradation on flat models (Gemma-2: unconstrained $1.47\times$ degradation reduced to within 0.5\% of uniform), while rendering solver choice irrelevant---constrained greedy and MCKP produce identical allocations---enabling a simple $O(H\log H)$ greedy algorithm.
We characterize sensitivity heterogeneity across all nine models using concentration and distribution shape ($t_1/t_5$ ratio), identifying three patterns (spread, spike, flat) that correctly predict allocation benefit (9/9); the thresholds are preliminary, as only one model exhibits the spread pattern.
